% GENERATED FILE. Edit canonical lessons, then run npm run generate:books.
% canonical-source-hash: fnv1a64:6055f95781e8581e
% canonical-lessons: GU-C281-dukh, GU-W50-visarga, GU-C281-mane-bik, GU-C281-iccha, GU-W51-short-i, GU-C281-sapnu, GU-C281-banvu-che, GU-R281-feelings-and-hopes

\chapter{Feelings and hopes}
\label{ch:gu-feelings-and-hopes}
\hlchaptermodality{281}

\begin{chapteropening}
\textbf{By the end of this chapter:} \emph{I can say how I feel and felt (mane bīk lāge che, hũ khuś chũ, mane duḥkh thayũ), say what I want to do and become (māre dāktar banvũ che), talk about a wish and a dream, and read and write \gu{ઃ} and \gu{ઇ}.}

Complete GU-R281-feelings-and-hopes: say how you feel, and what you want and hope for.
\end{chapteropening}

\section[\texorpdfstring{\gu{દુઃખ} (duḥkh)}{દુઃખ (duḥkh)}]{\gu{દુઃખ} (duḥkh) — sorrow; distress}

\label{lesson:GU-C281-dukh}



\begin{quote}
Before the new one: say \textquotedblleft{}contentment\textquotedblright{} and \textquotedblleft{}happiness\textquotedblright{}. (\textbf{\gu{સુખ}}, \textbf{\gu{ખુશી}}.) Then \textquotedblleft{}it happened\textquotedblright{}. (\textbf{\gu{થયું}}.)
\end{quote}

\subsection*{You'll want to know: \gu{દુઃખ}}
\textbf{\gu{દુઃખ}} — \emph{duḥkh} — \textquotedblleft{}sorrow; distress\textquotedblright{}, the opposite of \textbf{\gu{સુખ}}. It is neuter. In speech its \textbf{\gu{ઃ}} is light, close to \emph{dukh}.

\begin{quote}
\textbf{\gu{મને} \gu{દુઃખ} \gu{થયું}.}
\end{quote}

\begin{quote}
\emph{mane duḥkh thayũ.} — \textquotedblleft{}I felt sad.\textquotedblright{}
\end{quote}

Word for word, \textquotedblleft{}to me sorrow happened\textquotedblright{}: the one who feels it takes \textbf{-\gu{ને}}.

\begin{quote}
\textbf{\gu{તેને} \gu{કોઈ} \gu{દુઃખ} \gu{નથી}.}
\end{quote}

\begin{quote}
\emph{tene koī duḥkh nathī.} — \textquotedblleft{}He has no troubles.\textquotedblright{}
\end{quote}

\begin{cousinweb}
\textbf{\gu{દુઃખ}} (\emph{duḥkh}) is a tatsama: Sanskrit \emph{duḥkha}, \textquotedblleft{}difficulty, pain\textquotedblright{} (CDIAL 6375), written with the Sanskrit \textbf{\gu{ઃ}} kept. CDIAL also lists an inherited Gujarati \emph{dukh} from the same word.
\end{cousinweb}

\subsection*{Your turn}
Say these aloud:

\begin{itemize}
\raggedright
  \item \emph{duḥkh}, then \emph{sukh}
  \item \textquotedblleft{}I felt sad\textquotedblright{} — \emph{mane duḥkh thayũ}
  \item \textquotedblleft{}He has no troubles\textquotedblright{} — \emph{tene koī duḥkh nathī}
\end{itemize}

\begin{tcolorbox}[breakable,colback=teal!4,colframe=teal!35!black,title={Before you move on}]
What does \textbf{\gu{દુઃખ}} mean? (\textquotedblleft{}Sorrow; distress\textquotedblright{}.) Its opposite? (\textbf{\gu{સુખ}}.) \textquotedblleft{}I felt sad\textquotedblright{}? (\textbf{\gu{મને} \gu{દુઃખ} \gu{થયું}}.)
\end{tcolorbox}
\section[\texorpdfstring{\gu{ઃ} (visarga)}{ઃ (visarga)}]{\gu{ઃ} (visarga) — the two dots in \gu{દુઃખ}}

\label{lesson:GU-W50-visarga}



\begin{quote}
Before the new one: say \textbf{\gu{દુઃખ}} (\emph{duḥkh}), \textquotedblleft{}sorrow\textquotedblright{}. Then write \textbf{\gu{ં}}, the one dot above a syllable.
\end{quote}

\subsection*{Script you'll notice: \gu{ઃ}}
\textbf{\gu{ઃ}} — the visarga, \emph{ḥ}: a pair of dots to the right of the syllable, marking a voiceless breath after the vowel.

In \textbf{\gu{દુઃખ}} it stands after \textbf{\gu{દુ}} and before \textbf{\gu{ખ}}:

\noindent
\begin{tabularx}{\linewidth}{@{}>{\raggedright\arraybackslash}X>{\raggedright\arraybackslash}X@{}}
\toprule
\textbf{} & \textbf{} \\
\midrule
\textbf{\gu{દુ}} & \emph{du} \\
\textbf{\gu{દુઃ}} & \emph{duḥ}: the two dots after it \\
\textbf{\gu{દુઃખ}} & \emph{duḥkh} \\
\bottomrule
\end{tabularx}

It is written after its consonant and its vowel sign: the consonant first, then the visarga to its right, the upper dot, then the lower dot.

\subsection*{Writing — notice and trace}
\hlblockfigure{\detokenize{figures/GU-W50-visarga-filmstrip.pdf}}{How \gu{ઃ} is written, stroke by stroke}

Follow the numbered strip:

\begin{enumerate}
\raggedright
  \item start at the top of the upper dot and curve down its left side
  \item without lifting, round its bottom and up its right side to close it
  \item lift, then start at the top of the lower dot and curve down its left side
  \item without lifting, round its bottom and up its right side to close it — and only now lift
\end{enumerate}

Two strokes — four movements, one pen lift; cited to KanoAI's hand-made Gujarati barakhadi templates. Finger-trace \textbf{\gu{ઃ}} once.

\subsection*{Writing — copy with the model}
Copy \textbf{\gu{દુઃ}} once: \textbf{\gu{દુ}} first, then the two dots. Compare only two things: do the dots sit to the right, one above the other? Repair one if needed.

\subsection*{Writing — recall after hiding}
Hide the model and write the visarga once. Uncover \textbf{\gu{ઃ}}, compare, and stop after one repair.

\begin{tcolorbox}[breakable,colback=teal!4,colframe=teal!35!black,title={Before you move on}]
Read \textbf{\gu{દુઃખ}} aloud. Which dot comes first? (\textbf{The upper one.}) Is it written before or after \textbf{\gu{દુ}}? (After.)

Sources: \href{https://github.com/gajjartejas/KanoAI/tree/9d3e2949a3f265448e430e15329093ea3aa516b8/interpolate-svg/svgs/barakhadi}{KanoAI's Gujarati barakhadi stroke templates} (Tejas Gajjar, KanoAI, Gujarati barakhadi centre-line stroke templates (interpolate-svg/svgs/barakhadi), commit 9d3e294, 1\_k/11\_kah.svg, groups g1 then g2 (strokes g1s0, g2s0): the visarga drawn after the consonant (GitHub, 2026)) and the \href{https://www.unicode.org/charts/PDF/U0A80.pdf}{Unicode Gujarati chart}.
\end{tcolorbox}
\section[\texorpdfstring{\gu{મને} \gu{બીક} \gu{લાગે} \gu{છે} (mane bīk lāge che)}{મને બીક લાગે છે (mane bīk lāge che)}]{\gu{મને} \gu{બીક} \gu{લાગે} \gu{છે} (mane bīk lāge che) — how I feel}

\label{lesson:GU-C281-mane-bik}



\begin{quote}
Say \textquotedblleft{}I felt sad\textquotedblright{}. (\emph{mane duḥkh thayũ}.) Then \textquotedblleft{}fear\textquotedblright{} and \textquotedblleft{}worry\textquotedblright{}. (\textbf{\gu{બીક}}, \textbf{\gu{ચિંતા}}.)
\end{quote}

\begin{grammarlens}[title={Grammar Lens: a feeling that comes to me}]
\begin{quote}
\textbf{\gu{મને} \gu{બીક} \gu{લાગે} \gu{છે}.}
\end{quote}

\begin{quote}
\emph{mane bīk lāge che.} — \textquotedblleft{}I'm afraid.\textquotedblright{}
\end{quote}

Many feelings are nouns that happen to the one who feels them, and that one takes \textbf{-\gu{ને}}. Each noun has its own verb:

\noindent
\begin{tabularx}{\linewidth}{@{}>{\raggedright\arraybackslash}X>{\raggedright\arraybackslash}X>{\raggedright\arraybackslash}X@{}}
\toprule
\textbf{the feeling} & \textbf{now} & \textbf{in the past} \\
\midrule
fear, \textbf{\gu{બીક}} & \textbf{\gu{મને} \gu{બીક} \gu{લાગે} \gu{છે}} & \textbf{\gu{મને} \gu{બીક} \gu{લાગી}} \\
worry, \textbf{\gu{ચિંતા}} & \textbf{\gu{મને} \gu{ચિંતા} \gu{થાય} \gu{છે}} & \textbf{\gu{મને} \gu{ચિંતા} \gu{થઈ}} \\
sorrow, \textbf{\gu{દુઃખ}} & \textbf{\gu{મને} \gu{દુઃખ} \gu{થાય} \gu{છે}} & \textbf{\gu{મને} \gu{દુઃખ} \gu{થયું}} \\
anger, \textbf{\gu{ગુસ્સો}} & \textbf{\gu{મને} \gu{ગુસ્સો} \gu{આવે} \gu{છે}} & \textbf{\gu{મને} \gu{ગુસ્સો} \gu{આવ્યો}} \\
\bottomrule
\end{tabularx}

The past follows the feeling: \textbf{\gu{બીક}} is feminine, \textbf{\gu{લાગી}}; \textbf{\gu{ગુસ્સો}} is masculine, \textbf{\gu{આવ્યો}}.
\end{grammarlens}

\subsection*{You'll want to know: I am happy}
\begin{quote}
\textbf{\gu{આજે} \gu{હું} \gu{ખુશ} \gu{છું}.}
\end{quote}

\begin{quote}
\emph{āje hũ khuś chũ.} — \textquotedblleft{}Today I'm happy.\textquotedblright{}
\end{quote}

Other feelings are describing words, and the one who feels them is a plain subject: \textbf{\gu{હું} \gu{ખુશ} \gu{છું}}, \textbf{\gu{તે} \gu{ઉદાસ} \gu{હતી}} (\emph{te udās hatī}), \textquotedblleft{}she was sad\textquotedblright{}. Tiredness is a past: \textbf{\gu{હું} \gu{થાકી} \gu{ગયો} \gu{છું}} (\emph{hũ thākī gayo chũ}), \textquotedblleft{}I'm tired\textquotedblright{}, word for word \textquotedblleft{}I have got tired\textquotedblright{}.

\subsection*{Your turn}
Say these aloud:

\begin{itemize}
\raggedright
  \item \textquotedblleft{}I got angry\textquotedblright{} — \emph{mane gusso āvyo}
  \item \textquotedblleft{}My train was late, so I was worried\textquotedblright{} — \emph{mārī ṭren moḍī hatī, eṭle mane cintā thaī}
\end{itemize}

\begin{tcolorbox}[breakable,colback=teal!4,colframe=teal!35!black,title={Before you move on}]
\textquotedblleft{}I'm afraid\textquotedblright{}? (\textbf{\gu{મને} \gu{બીક} \gu{લાગે} \gu{છે}}.) \textquotedblleft{}I got angry\textquotedblright{}? (\textbf{\gu{મને} \gu{ગુસ્સો} \gu{આવ્યો}}.) \textquotedblleft{}She was sad\textquotedblright{}: \textbf{\gu{તે}} or \textbf{\gu{તેને}}? (\textbf{\gu{તે} \gu{ઉદાસ} \gu{હતી}}.)
\end{tcolorbox}
\section[\texorpdfstring{\gu{ઇચ્છા} (icchā)}{ઇચ્છા (icchā)}]{\gu{ઇચ્છા} (icchā) — a wish; a desire}

\label{lesson:GU-C281-iccha}



\begin{quote}
Before the new one: say \textquotedblleft{}I'm afraid\textquotedblright{}. (\emph{mane bīk lāge che}.) Then \textquotedblleft{}good wishes\textquotedblright{}. (\textbf{\gu{શુભેચ્છા}}, \emph{śubhecchā}.) Its second half is this lesson's word.
\end{quote}

\subsection*{You'll want to know: \gu{ઇચ્છા}}
\textbf{\gu{ઇચ્છા}} — \emph{icchā} — \textquotedblleft{}a wish; a desire\textquotedblright{}. It is feminine: \textbf{\gu{મારી} \gu{ઇચ્છા}}, \textquotedblleft{}my wish\textquotedblright{}.

\begin{quote}
\textbf{\gu{મારી} \gu{એક} \gu{ઇચ્છા} \gu{છે}.}
\end{quote}

\begin{quote}
\emph{mārī ek icchā che.} — \textquotedblleft{}I have a wish.\textquotedblright{}
\end{quote}

\begin{quote}
\textbf{\gu{તમારી} \gu{ઇચ્છા} \gu{શું} \gu{છે}?}
\end{quote}

\begin{quote}
\emph{tamārī icchā śũ che?} — \textquotedblleft{}What is your wish?\textquotedblright{}
\end{quote}

\textbf{\gu{શુભેચ્છા}} is \emph{śubh}, \textquotedblleft{}good, lucky\textquotedblright{}, with \textbf{\gu{ઇચ્છા}}: good wishes.

\begin{cousinweb}
\textbf{\gu{ઇચ્છા}} (\emph{icchā}) is a tatsama: Sanskrit \emph{icchā}, \textquotedblleft{}wish, desire\textquotedblright{} (Monier-Williams). CDIAL lists an inherited Gujarati verb \emph{ichvũ}, \textquotedblleft{}to wish\textquotedblright{}, under Sanskrit \emph{icchati}, \textquotedblleft{}seeks, wishes\textquotedblright{} (CDIAL 1556).
\end{cousinweb}

\subsection*{Your turn}
Say these aloud:

\begin{itemize}
\raggedright
  \item \emph{icchā}, then \emph{mārī icchā}
  \item \textquotedblleft{}I have a wish\textquotedblright{} — \emph{mārī ek icchā che}
  \item \textquotedblleft{}What is your wish?\textquotedblright{} — \emph{tamārī icchā śũ che?}
\end{itemize}

\begin{tcolorbox}[breakable,colback=teal!4,colframe=teal!35!black,title={Before you move on}]
What does \textbf{\gu{ઇચ્છા}} mean? (\textquotedblleft{}A wish; a desire\textquotedblright{}.) Why \textbf{\gu{મારી}}? (\textbf{\gu{ઇચ્છા}} is feminine.)
\end{tcolorbox}
\section[\texorpdfstring{\gu{ઇ} (i)}{ઇ (i)}]{\gu{ઇ} — the standing vowel i}

\label{lesson:GU-W51-short-i}



\begin{quote}
\begin{itemize}
\raggedright
  \item \emph{Recall:} say \emph{a wish} — \textbf{R1}, one lesson back
  \item \emph{Recall:} say \textbf{\gu{ઈ}}, as in \textbf{\gu{ઈજા}}, and the sign \textbf{\gu{િ}} that gives a consonant a short \emph{i}
\end{itemize}

\textbf{\gu{ઇચ્છા}} opens on a short \emph{i} with no consonant to carry it.
\end{quote}

\subsection*{You'll want to know: \gu{ઈ} without its curl}
\noindent
\begin{tabularx}{\linewidth}{@{}>{\raggedright\arraybackslash}X>{\raggedright\arraybackslash}X>{\raggedright\arraybackslash}X@{}}
\toprule
\textbf{} & \textbf{after a consonant} & \textbf{standing alone} \\
\midrule
short \emph{i} & \textbf{\gu{િ}} & ? \\
long \emph{ī} & \textbf{\gu{ી}} & \textbf{\gu{ઈ}} \\
\bottomrule
\end{tabularx}

The missing letter is \textbf{\gu{ઇ}}. It has the same small upper-left loop, narrow middle crossing and broad lower loop as \textbf{\gu{ઈ}}. Its hook rises on the upper right and stops there; \textbf{\gu{ઈ}} carries it on, curling over the top.

\subsection*{Script — meet one form}
\begin{quote}
\textbf{\gu{ઇ}} — \emph{i}
\end{quote}

Keep \textbf{\gu{ઇ}} visible. Follow the numbered strip:

\begin{enumerate}
\raggedright
  \item begin in the upper-left and circle the small loop down to the middle crossing
  \item continue through the crossing and sweep clockwise around the broad lower loop
  \item rise along the right side and finish in the upper-right hook without lifting
\end{enumerate}

One stroke, no lift: a verified one-stroke teaching-app variant fitted to the bundled printed outline (t30apps.com, ‘Gujarati Alphabet Writing Practice’, version 1.0, \gu{ઇ} animation, first SVG path).

\subsection*{Writing — notice and trace}
\hlblockfigure{\detokenize{figures/GU-W51-short-i-filmstrip.pdf}}{How \gu{ઇ} is written, stroke by stroke}

Finger-trace \textbf{\gu{ઇ}} once, then air-trace it once. Notice where the hook stops.

\subsection*{Writing — copy with the model}
Copy visible \textbf{\gu{ઇ}} once. Compare only two things: are both loops there, and does the hook stop on the upper right? Repair one landmark if needed.

\subsection*{Writing — recall after hiding}
Hide the model and write the standing vowel \emph{i} once. Uncover \textbf{\gu{ઇ}}, compare, and stop after one repair.

\begin{tcolorbox}[breakable,colback=teal!4,colframe=teal!35!black,title={Before you move on}]
Read \textbf{\gu{ઇ} \gu{ઈ} \gu{ઇચ્છા}} aloud. What turns \textbf{\gu{ઇ}} into \textbf{\gu{ઈ}}? (\textbf{The curl over the top.})

Sources: \href{https://www.t30apps.com/gujarati-alphabet-writing/}{Gujarati Alphabet Writing Practice} and the \href{https://www.unicode.org/charts/PDF/U0A80.pdf}{Unicode Gujarati chart}.
\end{tcolorbox}
\section[\texorpdfstring{\gu{સપનું} (sapnũ)}{સપનું (sapnũ)}]{\gu{સપનું} (sapnũ) — a dream; an ambition}

\label{lesson:GU-C281-sapnu}



\begin{quote}
Before the new one: say \textquotedblleft{}my wish\textquotedblright{}. (\emph{mārī icchā}.) Then \textquotedblleft{}sleepiness\textquotedblright{}. (\textbf{\gu{ઊંઘ}}.)
\end{quote}

\subsection*{You'll want to know: \gu{સપનું}}
\textbf{\gu{સપનું}} — \emph{sapnũ} — \textquotedblleft{}a dream; an ambition\textquotedblright{}. It is neuter: \textbf{\gu{મારું} \gu{સપનું}}. In Gujarati a dream comes to the dreamer:

\begin{quote}
\textbf{\gu{મને} \gu{રાતે} \gu{એક} \gu{સપનું} \gu{આવ્યું}.}
\end{quote}

\begin{quote}
\emph{mane rāte ek sapnũ āvyũ.} — \textquotedblleft{}I had a dream last night.\textquotedblright{}
\end{quote}

The second sense is a hope for the future:

\begin{quote}
\textbf{\gu{મારું} \gu{સપનું} \gu{છે} \gu{કે} \gu{એક} \gu{દિવસ} \gu{હું} \gu{દાક્તર} \gu{બનીશ}.}
\end{quote}

\begin{quote}
\emph{mārũ sapnũ che ke ek divas hũ dāktar banīś.} — \textquotedblleft{}My dream is that one day I'll become a doctor.\textquotedblright{}
\end{quote}

\begin{cousinweb}
\textbf{\gu{સપનું}} (\emph{sapnũ}) is usually traced to Sanskrit \emph{svapna}, \textquotedblleft{}sleep; a dream\textquotedblright{} (Monier-Williams), worn down in speech: the \emph{sv-} became \emph{s-}. CDIAL lists a related Gujarati \emph{samṇũ}, \textquotedblleft{}dream\textquotedblright{}, under a reconstructed \emph{supna} (CDIAL 13481).
\end{cousinweb}

\subsection*{Your turn}
Say these aloud:

\begin{itemize}
\raggedright
  \item \emph{sapnũ}, then \emph{mārũ sapnũ}
  \item \textquotedblleft{}I had a dream last night\textquotedblright{} — \emph{mane rāte ek sapnũ āvyũ}
  \item \textquotedblleft{}My dream is that one day I'll become a doctor\textquotedblright{} — \emph{mārũ sapnũ che ke ek divas hũ dāktar banīś}
\end{itemize}

\begin{tcolorbox}[breakable,colback=teal!4,colframe=teal!35!black,title={Before you move on}]
What does \textbf{\gu{સપનું}} mean? (\textquotedblleft{}A dream; an ambition\textquotedblright{}.) \textquotedblleft{}I had a dream\textquotedblright{}? (\textbf{\gu{મને} \gu{સપનું} \gu{આવ્યું}}.)
\end{tcolorbox}
\section[\texorpdfstring{\gu{મારે} \gu{દાક્તર} \gu{બનવું} \gu{છે} (māre dāktar banvũ che)}{મારે દાક્તર બનવું છે (māre dāktar banvũ che)}]{\gu{મારે} \gu{દાક્તર} \gu{બનવું} \gu{છે} (māre dāktar banvũ che) — what I want, and hope for}

\label{lesson:GU-C281-banvu-che}



\begin{quote}
Say \textquotedblleft{}I should go\textquotedblright{}. (\emph{māre javũ joīe}.) Then \textquotedblleft{}my dream\textquotedblright{} and \textquotedblleft{}my wish\textquotedblright{}. (\emph{mārũ sapnũ}, \emph{mārī icchā}.)
\end{quote}

\begin{grammarlens}[title={Grammar Lens: \gu{મારે} … -\gu{વું} \gu{છે}}]
\begin{quote}
\textbf{\gu{મારે} \gu{દાક્તર} \gu{બનવું} \gu{છે}.}
\end{quote}

\begin{quote}
\emph{māre dāktar banvũ che.} — \textquotedblleft{}I want to become a doctor.\textquotedblright{}
\end{quote}

The \textbf{\gu{મારે}} of \textbf{\gu{મારે} \gu{જવું} \gu{જોઈએ}}, then the \textbf{-\gu{વું}} form, then \textbf{\gu{છે}}: \textquotedblleft{}I want to\textquotedblright{}. With an object, the \textbf{-\gu{વું}} follows it, as with \textbf{\gu{જોઈએ}}: \textbf{\gu{મારે} \gu{ચા} \gu{પીવી} \gu{છે}} (\emph{māre cā pīvī che}), \textquotedblleft{}I want to drink tea\textquotedblright{}; \textbf{\gu{મારે} \gu{નવું} \gu{ઘર} \gu{બનાવવું} \gu{છે}}. \textquotedblleft{}I don't want to go\textquotedblright{} is \textbf{\gu{મારે} \gu{નથી} \gu{જવું}} (\emph{māre nathī javũ}).
\end{grammarlens}

\subsection*{You'll want to know: a wish, a dream}
\begin{quote}
\textbf{\gu{મારી} \gu{ઇચ્છા} \gu{ગામમાં} \gu{રહેવાની} \gu{છે}.}
\end{quote}

\begin{quote}
\emph{mārī icchā gāmmā̃ rahevānī che.} — \textquotedblleft{}My wish is to live in the village.\textquotedblright{}
\end{quote}

\textbf{\gu{રહેવાની}} is the \textbf{-\gu{વાનો}} of \textbf{\gu{હું} \gu{જવાનો} \gu{છું}}, here following \textbf{\gu{ઇચ્છા}}, which is feminine. With \textbf{\gu{સપનું}}, neuter, it is \textbf{-\gu{વાનું}}: \textbf{\gu{મારું} \gu{સપનું} \gu{દાક્તર} \gu{બનવાનું} \gu{છે}} (\emph{mārũ sapnũ dāktar banvānũ che}). A dream can also take \textbf{\gu{કે}} and a future: \textbf{\gu{મારું} \gu{સપનું} \gu{છે} \gu{કે} \gu{હું} \gu{દાક્તર} \gu{બનીશ}}.

\subsection*{Your turn}
Say these aloud:

\begin{itemize}
\raggedright
  \item \textquotedblleft{}I want to drink tea\textquotedblright{} — \emph{māre cā pīvī che}
  \item \textquotedblleft{}I don't want to go to the city\textquotedblright{} — \emph{māre śaher nathī javũ}
\end{itemize}

\begin{tcolorbox}[breakable,colback=teal!4,colframe=teal!35!black,title={Before you move on}]
\textquotedblleft{}I want to become a doctor\textquotedblright{}? (\textbf{\gu{મારે} \gu{દાક્તર} \gu{બનવું} \gu{છે}}.) Why \textbf{\gu{પીવી}} in \textbf{\gu{મારે} \gu{ચા} \gu{પીવી} \gu{છે}}? (It follows \textbf{\gu{ચા}}.) \textquotedblleft{}My dream is to become a doctor\textquotedblright{}? (\textbf{\gu{મારું} \gu{સપનું} \gu{દાક્તર} \gu{બનવાનું} \gu{છે}}.)
\end{tcolorbox}
\section[(review)]{Review — feelings and hopes}

\label{lesson:GU-R281-feelings-and-hopes}



\begin{quote}
Say \textquotedblleft{}my bag is lost\textquotedblright{}. (\emph{mārī thelī khovāī gaī}.) Then \textquotedblleft{}I'll have to go\textquotedblright{}. (\emph{māre javũ paḍśe}.)

\begin{itemize}
\raggedright
  \item \emph{Write it:} \textbf{\gu{ઃ}} in \textbf{\gu{દુઃખ}}, then \textbf{\gu{ઇ}} in \textbf{\gu{ઇચ્છા}}, from memory
\end{itemize}
\end{quote}

\subsection*{Your turn}
Say these aloud:

\begin{itemize}
\raggedright
  \item I'm afraid. I got angry. I felt sad.
  \item Today I'm happy. She was sad. I'm tired.
  \item I want to become a doctor. I don't want to go to the city.
  \item My wish is to live in the village.
  \item I had a dream last night. My dream is that one day I'll become a doctor.
\end{itemize}

\begin{tcolorbox}[breakable,colback=teal!4,colframe=teal!35!black,title={Before you move on}]
\textquotedblleft{}I'm afraid\textquotedblright{}: \textbf{\gu{હું}} or \textbf{\gu{મને}}? (\textbf{\gu{મને} \gu{બીક} \gu{લાગે} \gu{છે}}.) \textquotedblleft{}I'm happy\textquotedblright{}? (\textbf{\gu{હું} \gu{ખુશ} \gu{છું}}.) \textquotedblleft{}I want to drink tea\textquotedblright{}? (\textbf{\gu{મારે} \gu{ચા} \gu{પીવી} \gu{છે}}.) Which two dots sit in \textbf{\gu{દુઃખ}}? (\textbf{\gu{ઃ}}, the visarga.)
\end{tcolorbox}
